\section{Muons}
\label{sec:ch6_muons}

\subsection{Reconstruction}

Muons are reconstructed using tracks in the ID and MS, together with their energy deposits in the calorimeters~\cite{ATLAS2021MuonID}.
The main reconstruction categories are:
\begin{itemize}
    \item \textbf{Combined muons} are reconstructed by matching and fitting ID and MS tracks, including the energy loss between the two tracking systems.
    \item \textbf{Segment-tagged muons} are reconstructed by matching an ID track to one or more local track segments in the MS when a complete MS track is unavailable.
    \item \textbf{Calorimeter-tagged muons} are reconstructed by matching an ID track to calorimeter deposits compatible with a minimum-ionising particle. Acceptance is recovered in regions with limited MS coverage.
    \item \textbf{MS-extrapolated muons} are reconstructed by extrapolating an MS track towards the interaction region without a matched ID track.
\end{itemize}

Combined muons are used in this analysis.

\subsection{Identification}

Muon identification tests whether the detector hits and the ID and MS momentum measurements are consistent with a single muon trajectory~\cite{ATLAS2021MuonID}.
Decays of charged mesons in flight can change the trajectory and momentum, leading to a disagreement between the momentum measurements.
For muon candidates, this agreement is quantified by the $q/p$ significance,
\begin{align}
    S_{q/p}=\frac{\left|(q/p)_{\mathrm{ID}}-(q/p)_{\mathrm{MS}}\right|}
    {\sqrt{\sigma_{\mathrm{ID}}^2+\sigma_{\mathrm{MS}}^2}},
    \label{eq:ch6_muon_qp_significance}
\end{align}
where $(q/p)_{\mathrm{ID}}$ and $(q/p)_{\mathrm{MS}}$ are the charge-to-momentum ratios determined from the ID and MS tracks, respectively.
$\sigma_{\mathrm{ID}}$ and $\sigma_{\mathrm{MS}}$ denote the uncertainties.
Large values indicate poor agreement and help reject incorrectly matched candidates.

In the analysis, baseline muons are required to pass \texttt{Loose} identification while signal muons must additionally satisfy the \texttt{Medium} working point.

\subsection{Isolation}

Muons from heavy-flavour decays are often accompanied by other particles inside a jet.
They are suppressed using track and calorimeter isolation, defined by the surrounding activity after excluding the muon's own contribution~\cite{ATLAS2021MuonID}.
The \texttt{Loose\_VarRad} and \texttt{Tight\_VarRad} working points combine track and calorimeter isolation with a track-cone radius that decreases as the muon $\pT$ increases.
Less surrounding activity is allowed by \texttt{Tight\_VarRad}, giving stronger rejection of muons produced inside jets.
Baseline muons are required to pass \texttt{Loose\_VarRad} isolation, and signal muons must additionally satisfy \texttt{Tight\_VarRad}.

\subsection{Calibration}

The muon momentum scale and resolution are calibrated using resonance decays, particularly $\zmumu$ and $J/\psi\rightarrow\mu\mu$~\cite{ATLAS2023MuonCalibration}.
The reconstructed mass distributions constrain biases from the magnetic-field description, detector alignment and material modelling.
The corrections are applied to the combined momentum measurement used in the analysis.
Reconstruction, identification and isolation scale factors are measured with tag-and-probe samples and applied to simulated event weights for the selected working points~\cite{ATLAS2021MuonID}.
